\documentclass[11pt]{article}
\usepackage[margin=1in]{geometry}
\usepackage{amsmath,amssymb,amsthm,hyperref}
\newtheorem{theorem}{Theorem}
\newtheorem{lemma}{Lemma}
\newcommand{\Sym}{\operatorname{Sym}}
\newcommand{\sgn}{\operatorname{sgn}}
\newcommand{\ord}{\operatorname{ord}}
\title{Vanishing in Nuttall's normalized determinant expansion\\
\large A solution of MathDB 322495}
\author{Principia Math In Progress}
\date{27 September 2026}
\begin{document}
\maketitle
\begin{abstract}
We prove Conjecture 3 of Nuttall's \emph{Determinantal approach to a proof
of the Riemann hypothesis}. In fact, its coefficients vanish whenever
$m+2i<r(r-1)$, a slightly larger range than requested. The proof is purely
algebraic: the coefficient of the normalization parameter is a symmetric
differential operator of controlled order applied to a squared Vandermonde.
\end{abstract}

\section*{Task 1: the statement and its normalization}
Fix an integer $r>2$ and put $R=r(r-1)$. In \cite[Section 7.2]{N},
the coefficients $C(i;m,j)$ arise from the permutation expansion of
\[
 D(n,r)=\det_{1\leq a,b\leq r}
 \left[\frac{b_{n+b-a}}{\Gamma(2(n+b-a)+1)}\right],
 \qquad y=\frac1{2n},\quad n>r.
\]
Here $b_h$ are moments; their values play no role in the coefficient
identity. The normalization is
\[
 \overline D(n,r)=y^R\Gamma(2n+1)^r
 \prod_{h=1}^{r-1}\bigl((2n+2h)(2n+2h-1)\bigr)^{r-h}D(n,r).
\]
Nuttall asks for $C(i;m,j)=0$ when
\[
 0\leq i\leq R/2-1,\qquad 0\leq m\leq R-2i-2,
 \qquad 1\leq j\leq np(m).
\]
We specify these coefficients without integrals. Symmetrization below is
the average over all permutations of the $r$ variables. For
$\sigma\in S_r$ define
\begin{align}
 \nu_a&=\sigma(a)-a,\qquad
 M_\sigma(t)=\prod_{a=1}^r t_a^{2(\nu_a+r-1)},\label{eq:monomial}\\
 Z_\sigma(y)&=\prod_{a=1}^r
       \prod_{s=2(\sigma(a)-a)+1}^{2(r-a)}(1+sy).\label{eq:weight}
\end{align}
An empty product is $1$. All exponents in \eqref{eq:monomial} are
nonnegative, and every $Z_\sigma$ is a polynomial of degree at most $R$.
Set
\begin{equation}\label{eq:P}
 P_r(y,t)=\sum_{\sigma\in S_r}\sgn(\sigma)Z_\sigma(y)
                                  \Sym M_\sigma(t).
\end{equation}
To check \eqref{eq:weight}, write $z=2n$. Multiplying row $a$ of the
gamma-ratio matrix by $\Gamma(z+2(r-a)+1)/\Gamma(z+1)$ and its
column $b$ by $y^{2(r-b)}$ gives the entry
\[
 y^{2(r-b)}\frac{\Gamma(z+2(r-a)+1)}{\Gamma(z+2(b-a)+1)}
 =\prod_{s=2(b-a)+1}^{2(r-a)}(1+sy).
\]
The product of these row and column factors is precisely the stated
normalization of the gamma-ratio determinant. Thus $Z_\sigma$ is its
unsigned permutation component.

Write $t_a=1+x_a$. The source's symmetrized monomials
$\mu(m,j)$ form a basis for symmetric homogeneous polynomials of degree
$m$. Its equations (7.17)--(7.22) therefore say exactly that
\begin{equation}\label{eq:C}
 P_r(y,1+x)=\sum_{i,m,j}C(i;m,j)y^i\mu(m,j)(x).
\end{equation}
Indeed, replacing the source variables by $\tau(1+x_a)$ extracts the
common factor $\tau^{2R}$ from every $M_\sigma$, as in (7.19).
This also fixes all normalizations of $C$.

\begin{theorem}\label{thm:main}
For every $r\geq1$, the coefficient of $y^i$ in $P_r(y,1+x)$ has no
monomial of total $x$-degree less than $R-2i$. Consequently
\[
 C(i;m,j)=0\quad\hbox{whenever}\quad m+2i<R.
\]
In particular, all the coefficients required by Conjecture 3 vanish.
\end{theorem}

\section*{Polynomial dependence of the weights}
For any integer $v$ define a formal power series at $y=0$ by
\[
 F_v(y)=
 \begin{cases}
 \displaystyle\prod_{s=1}^{2v}(1+sy)^{-1},&v\geq0,\\[4pt]
 \displaystyle\prod_{s=2v+1}^{0}(1+sy),&v<0.
 \end{cases}
\]
Also put
\[
 K_r(y)=\prod_{h=1}^{r-1}\bigl((1+2hy)(1+(2h-1)y)\bigr)^{r-h}.
\]
Since $\sum_a\nu_a=0$, cancellation of the finite products, or the
gamma-ratio calculation above, gives
\begin{equation}\label{eq:factor}
 Z_\sigma(y)=K_r(y)\prod_{a=1}^r F_{\nu_a}(y).
\end{equation}

\begin{lemma}\label{lem:degree}
The coefficient $[y^k]F_v(y)$ is a polynomial in $v$ of degree at most
$2k$, valid for positive, zero, and negative integers $v$. Hence there
is a symmetric polynomial $Q_i(v_1,\ldots,v_r)$, of total degree at
most $2i$, such that
\[
 [y^i]Z_\sigma(y)=Q_i(\nu_1,\ldots,\nu_r).
\]
\end{lemma}
\begin{proof}
Let $S_\ell(u)$ be the polynomial of degree $\ell+1$ satisfying
$S_\ell(0)=0$ and
$S_\ell(u)-S_\ell(u-1)=u^\ell$. Such a polynomial exists because the
finite-difference operator takes degree $\ell+1$ to degree $\ell$ with
nonzero leading coefficient. Telescoping identifies it with
$\sum_{s=1}^{u}s^\ell$ for positive integers $u$, and with
$-\sum_{s=u+1}^{0}s^\ell$ for negative integers $u$. Thus, in both cases,
\[
 \log F_v(y)=\sum_{\ell\geq1}\frac{(-1)^\ell}{\ell}
                              S_\ell(2v)y^\ell.
\]
In the coefficient of $y^k$ in its exponential, each product has
$\sum_\ell \ell a_\ell=k$ and degree at most
$\sum_\ell(\ell+1)a_\ell\leq2k$ in $v$.
Multiplication over the $r$ variables preserves the total-degree bound
and gives a symmetric polynomial. The scalar series $K_r(y)$ preserves
both properties, proving the assertion.
\end{proof}

\section*{Symmetrization and differential order}
Define commuting operators
\[
 E_a=\frac12t_a\frac{\partial}{\partial t_a}-(r-1),
 \qquad 1\leq a\leq r.
\]
Equation \eqref{eq:monomial} gives $E_aM_\sigma=\nu_aM_\sigma$.
Since $Q_i$ is symmetric, $Q_i(E_1,\ldots,E_r)$ commutes with
symmetrization. Therefore
\begin{equation}\label{eq:operator}
 [y^i]P_r(y,t)=Q_i(E_1,\ldots,E_r)B_r(t),\qquad
 B_r(t)=\Sym\sum_{\sigma\in S_r}\sgn(\sigma)M_\sigma(t).
\end{equation}
The operator on the right has differential order at most $2i$.

Put $u_a=t_a^2$ and $V(u)=\prod_{a<b}(u_b-u_a)$. The determinant
expansion shows that
\[
 \sum_{\sigma\in S_r}\sgn(\sigma)M_\sigma(t)
 =\left(\prod_{a=1}^r u_a^{r-a}\right)V(u).
\]
When we symmetrize, $V(u)$ contributes the sign of the variable
permutation. The resulting alternating sum of the descending powers
$r-1,r-2,\ldots,0$ is $(-1)^{R/2}V(u)$. Hence
\begin{equation}\label{eq:V}
 B_r(t)=\frac{(-1)^{R/2}}{r!}\prod_{a<b}(t_b^2-t_a^2)^2.
\end{equation}

\begin{proof}[Proof of Theorem \ref{thm:main}]
At $t=1+x$ every factor in \eqref{eq:V} has the form
\[
 (t_b^2-t_a^2)^2=(x_b-x_a)^2(2+x_a+x_b)^2.
\]
Thus $B_r(1+x)$ belongs to the $R$th power of the ideal
$\mathfrak m=(x_1,\ldots,x_r)$. A differential operator of order at most
$d$ with polynomial coefficients maps $\mathfrak m^R$ into
$\mathfrak m^{\max(R-d,0)}$: each derivative lowers total degree by at
most one, and multiplication by a polynomial cannot lower it.
Apply this with $d=2i$ to \eqref{eq:operator}. All homogeneous components
of degree $m<R-2i$ vanish. Their coefficients in the independent basis
$\mu(m,j)$ in \eqref{eq:C} therefore vanish individually.
\end{proof}

The proof is valid for arbitrary formal moments and uses no analytic
assumption on Nuttall's function $\Phi$. Its conclusion is the algebraic
coefficient identity. The analytic positivity and asymptotic questions
motivating the source require additional arguments.

\section*{Verification and proposed follow-up tasks}
The accompanying standard-library Python verifier directly expands the
permutation products \eqref{eq:weight} and the symmetrized binomial
coefficients. It checks 2,810 zero coefficients in ranks 2--5, including
the extra degree covered by Theorem \ref{thm:main}. It separately compares
the product normalization with exact factorial ratios. These finite
checks supplement the all-rank proof.

\begin{description}
\item[Task 2 (proposed).] Give a closed formula for the part of $P_r(y,1+x)$
of weighted degree $R$, with weights $2$ for $y$ and $1$ for every $x_a$.
\item[Task 3 (proposed).] Determine the exact order of vanishing at $x=0$
of $[y^i]P_r(y,1+x)$ for every $0\leq i\leq R/2$, including when the
bound $R-2i$ is attained.
\item[Task 4 (proposed).] Replace the gamma arguments $2(n+b-a)+1$ by
$a_0(n+b-a)+b_0$, with integer $a_0>0$, and determine the optimal
coefficient-vanishing range under the corresponding moment normalization.
\item[Task 5 (proposed).] For Nuttall's particular $\Phi$, combine the
vanishing identity with a saddle-point expansion to determine the first
nonzero asymptotic term of $\overline D(n,r)$, with a proved remainder
bound for each fixed $r$.
\end{description}

\begin{thebibliography}{1}
\bibitem{N} J. Nuttall, \emph{Determinantal approach to a proof of the
Riemann hypothesis}, arXiv:1111.1128v1 (2011), Sections 6.5 and 7,
especially Conjecture 3 and equations (7.15)--(7.24).
\url{https://arxiv.org/abs/1111.1128}.
\end{thebibliography}
\smallskip
\noindent\textit{Review status.} Author-agent review; no independent peer
review or proof-assistant verification. Dated scope and claim checks
are documented in the accompanying audits.
\end{document}
